\documentclass[11pt]{article}
\usepackage{fontspec,amsmath,amssymb,amsthm,hyperref}
\usepackage[margin=1in]{geometry}
\setlength{\emergencystretch}{3em}
\title{Actual past and future Brownian independence and its L\ensuremath{{}^{2}} closedness}
\author{LangevinLean mathematical exposition}
\date{}
\begin{document}\maketitle
\noindent Source module: \texttt{KineticLangevin/PastFutureIndependence.lean}. Compile this proof document with XeLaTeX or LuaLaTeX. The accompanying HTML pairs the mathematical argument with the exact Lean statements and records the scope of the formalization.


The future convolution has support after the terminal time of every past kernel. Their actual isonormal Hilbert inner products vanish. Joint Gaussianity therefore proves independence of the whole past block and the future response. Independent spatial copies retain this property. The final argument proves independence survives genuine L\ensuremath{{}^{2}} limits, which lets numerical approximants establish independence for the actual nonlinear reference solution.


\section{Regroup independent arrays without changing their laws}


The exact independent array law is the corresponding infinite product. The measurable currying equivalence transports this law to the product of row laws. Pulling the independent product coordinate projections back along that actual law proves row independence. This handles the change from individual spatial coordinates to entire past and future spatial blocks.



\noindent Lean declarations: \texttt{iIndepFun\_curry}.

\section{Use disjoint deterministic time supports}


Place all past kernels in one Gaussian block and the future kernel in a second block. The joint family is Gaussian because every coordinate is the same actual isonormal linear map applied to its deterministic time kernel. A past interval ends before the future begins, so every cross covariance is exactly zero. Gaussian block independence follows. Independent coordinate source projections then assemble the full rectangular spatial array, and currying groups it into independent past and future vector blocks.



\noindent Lean declarations: \texttt{scalarSegments\_future\_independent}, \texttt{vectorSegments\_future\_blocks\_independent}.

\section{Recover the actual Euclidean convolution variables}


Pack each spatial row with the actual Euclidean coordinate map and read the future response from any nonempty block coordinate. The spatial-lift representative theorem identifies these packed functions with the genuine vector L\ensuremath{{}^{2}} segment convolutions. A countable almost-everywhere intersection handles the entire past family. Independence is unchanged by these almost-everywhere representative identities.



\noindent Lean declarations: \texttt{vectorSegments\_future\_independent}.

\section{Adjoin initial data and change probability sources}


Measure-preserving source composition transports independence by exact pushforward composition. For a separate independent initial coordinate, product mapping first gives the joint initial/past/future law. Product associativity then places initial data and past together opposite the future coordinate. The two product projections are independent, proving that arbitrary measurable functions of the initial data and past remain independent of the future.



\noindent Lean declarations: \texttt{indepFun\_comp\_measurePreserving}, \texttt{indepFun\_prod\_initial\_of\_indep}.

\section{Pass independence through actual L\ensuremath{{}^{2}} limits}


A continuous linear functional sends an L\ensuremath{{}^{2}} vector to an L\ensuremath{{}^{2}} scalar. Continuity of the actual scalar characteristic function therefore gives continuity of every dual characteristic function of the vector law. Independence is characterized by factorization of the joint dual characteristic function. For a convergent sequence of past L\ensuremath{{}^{2}} variables and a fixed future variable, both the joint characteristic function and the product of marginal characteristic functions converge. Their equality at every approximation passes to the limit, proving independence of the actual limiting past variable.



\noindent Lean declarations: \texttt{continuous\_charFunDual\_map\_Lp}, \texttt{indepFun\_of\_tendsto\_Lp}.

\end{document}
